\section{Bibliografía}

\begin{thebibliography}{99}

% ── Normativa colombiana ────────────────────────────────────────────────────

\bibitem{NSR10}
Asociación Colombiana de Ingeniería Sísmica (AIS). (2010).
\textit{Reglamento Colombiano de Construcción Sismo Resistente NSR-10, Título C: Concreto Estructural.}
Ministerio de Ambiente, Vivienda y Desarrollo Territorial.

\bibitem{NTC673}
Instituto Colombiano de Normas Técnicas y Certificación (ICONTEC). (2010).
\textit{NTC 673: Ensayo de resistencia a la compresión de especímenes cilíndricos de concreto.}
ICONTEC.

\bibitem{NTC1299}
Instituto Colombiano de Normas Técnicas y Certificación (ICONTEC). (2004).
\textit{NTC 1299: Aditivos químicos para concreto.}
ICONTEC.

% ── Normativa internacional ─────────────────────────────────────────────────

\bibitem{ACI211}
ACI Committee 211. (2016).
\textit{ACI 211.1-16: Standard Practice for Selecting Proportions for Normal, Heavyweight, and Mass Concrete.}
American Concrete Institute.

\bibitem{ASTM_C39}
ASTM International. (2021).
\textit{ASTM C39/C39M-21: Standard Test Method for Compressive Strength of Cylindrical Concrete Specimens.}
ASTM International.

\bibitem{ASTM_C192}
ASTM International. (2021).
\textit{ASTM C192/C192M-21: Standard Practice for Making and Curing Concrete Test Specimens in the Laboratory.}
ASTM International.

\bibitem{ASTM_C494}
ASTM International. (2022).
\textit{ASTM C494/C494M-22: Standard Specification for Chemical Admixtures for Concrete.}
ASTM International.

% ── Literatura académica ────────────────────────────────────────────────────

\bibitem{jimenez2020}
Jiménez, L., \& Alvarado, M. (2020).
Modelos matemáticos para la optimización de mezclas de concreto: una revisión.
\textit{Revista de Ingeniería Civil}, 15(2), 45--58.

\bibitem{garcia2021}
García, P., \& Torres, R. (2021).
Integración de materiales cementantes alternativos en el diseño de mezclas de concreto.
\textit{Journal of Sustainable Construction}, 8(3), 112--125.

\bibitem{torres2023}
Torres, L., \& Jiménez, F. (2023).
Aplicación de redes neuronales en la predicción de propiedades del concreto.
\textit{Inteligencia Artificial en la Construcción}, 5(2), 78--90.

\bibitem{gonzalez2019}
González, R., Patiño, H., \& Ramírez, C. (2019).
Herramientas digitales para el diseño de mezclas de concreto en el contexto colombiano.
\textit{Revista Ingeniería y Región}, 21(1), 33--44.

\bibitem{tashildar2020}
Tashildar, A., Shah, N., Gala, R., Giri, T., \& Chavhan, P. (2020).
Application development using Flutter.
\textit{International Research Journal of Modernization in Engineering Technology and Science}, 2(8), 1262--1266.

\bibitem{palmer2002}
Palmer, S. R., \& Felsing, J. M. (2002).
\textit{A Practical Guide to Feature-Driven Development.}
Prentice Hall.

% ── Política pública ────────────────────────────────────────────────────────

\bibitem{DNP2018}
Departamento Nacional de Planeación (DNP). (2018).
\textit{Política Nacional de Edificaciones Sostenibles} (Documento CONPES 3919).
DNP.

\bibitem{CCI2020}
Cámara Colombiana de la Infraestructura (CCI). (2020).
\textit{Informe de sostenibilidad e innovación tecnológica en el sector de la construcción en Colombia.}
CCI.

% ── Software y herramientas relacionadas ────────────────────────────────────

\bibitem{concrebrain}
ConcreBrain. (s.f.). \textit{ConcreBrain: Software de predicción de resistencia del concreto mediante redes neuronales.} Recuperado de \url{https://concrebrain.com}

\bibitem{concreteMixDesign}
Concrete Mix Design. (s.f.). \textit{Concrete Mix Design: Software comercial para diseño de mezclas de concreto.}

\bibitem{concretelab}
Concretelab. (s.f.). \textit{Concretelab: Herramienta web para diseño y análisis de mezclas de concreto.}

% ── Tecnología usada en DIAPCE ──────────────────────────────────────────────

\bibitem{flutterdev}
Google LLC. (2024). \textit{Flutter — Build apps for any screen.} Recuperado de \url{https://flutter.dev}

\bibitem{dartdev}
Google LLC. (2024). \textit{Dart Programming Language.} Recuperado de \url{https://dart.dev}

\bibitem{sqliteDocs}
SQLite Consortium. (2024). \textit{SQLite Documentation.} Recuperado de \url{https://www.sqlite.org/docs.html}

\bibitem{sqflite}
Tekartik. (2024). \textit{sqflite — SQLite plugin for Flutter.} Recuperado de \url{https://pub.dev/packages/sqflite}

\bibitem{githubdocs}
GitHub Docs. (s.f.). \textit{GitHub Documentation.} Recuperado de \url{https://docs.github.com}

\bibitem{vscodedocs}
Visual Studio Code Docs. (s.f.). \textit{Visual Studio Code Documentation.} Recuperado de \url{https://code.visualstudio.com/docs}

% ── Metodología ─────────────────────────────────────────────────────────────

\bibitem{lucidchart2024}
Lucidchart. (2024). \textit{Cómo utilizar el desarrollo basado en funciones.} Recuperado de \url{https://www.lucidchart.com/blog/es/utilizar-el-desarrollo-basado-en-funciones}

\bibitem{fddsite}
Feature Driven Development. (s.f.). \textit{Feature Driven Development.} Recuperado de \url{https://featuredrivendevelopment.com}

\end{thebibliography}